\IfFileExists{stacks-project.cls}{%
\documentclass{stacks-project}
}{%
\documentclass{amsart}
}

% For dealing with references we use the comment environment
\usepackage{verbatim}
\newenvironment{reference}{\comment}{\endcomment}
%\newenvironment{reference}{}{}
\newenvironment{slogan}{\comment}{\endcomment}
\newenvironment{history}{\comment}{\endcomment}

% For commutative diagrams we use Xy-pic
\usepackage[all]{xy}

% We use 2cell for 2-commutative diagrams.
\xyoption{2cell}
\UseAllTwocells

% We use multicol for the list of chapters between chapters
\usepackage{multicol}

% This is generally recommended for better output
\usepackage{lmodern}
\usepackage[T1]{fontenc}

% For cross-file-references
\usepackage{xr-hyper}

% Package for hypertext links:
\usepackage{hyperref}

% For any local file, say "hello.tex" you want to link to please
% use \externaldocument[hello-]{hello}
\externaldocument[introduction-]{introduction}
\externaldocument[conventions-]{conventions}
\externaldocument[sets-]{sets}
\externaldocument[categories-]{categories}
\externaldocument[topology-]{topology}
\externaldocument[sheaves-]{sheaves}
\externaldocument[sites-]{sites}
\externaldocument[stacks-]{stacks}
\externaldocument[fields-]{fields}
\externaldocument[algebra-]{algebra}
\externaldocument[brauer-]{brauer}
\externaldocument[homology-]{homology}
\externaldocument[derived-]{derived}
\externaldocument[simplicial-]{simplicial}
\externaldocument[more-algebra-]{more-algebra}
\externaldocument[smoothing-]{smoothing}
\externaldocument[modules-]{modules}
\externaldocument[sites-modules-]{sites-modules}
\externaldocument[injectives-]{injectives}
\externaldocument[cohomology-]{cohomology}
\externaldocument[sites-cohomology-]{sites-cohomology}
\externaldocument[dga-]{dga}
\externaldocument[dpa-]{dpa}
\externaldocument[sdga-]{sdga}
\externaldocument[hypercovering-]{hypercovering}
\externaldocument[schemes-]{schemes}
\externaldocument[constructions-]{constructions}
\externaldocument[properties-]{properties}
\externaldocument[morphisms-]{morphisms}
\externaldocument[coherent-]{coherent}
\externaldocument[divisors-]{divisors}
\externaldocument[limits-]{limits}
\externaldocument[varieties-]{varieties}
\externaldocument[topologies-]{topologies}
\externaldocument[descent-]{descent}
\externaldocument[perfect-]{perfect}
\externaldocument[more-morphisms-]{more-morphisms}
\externaldocument[flat-]{flat}
\externaldocument[groupoids-]{groupoids}
\externaldocument[more-groupoids-]{more-groupoids}
\externaldocument[etale-]{etale}
\externaldocument[chow-]{chow}
\externaldocument[intersection-]{intersection}
\externaldocument[pic-]{pic}
\externaldocument[weil-]{weil}
\externaldocument[adequate-]{adequate}
\externaldocument[dualizing-]{dualizing}
\externaldocument[duality-]{duality}
\externaldocument[discriminant-]{discriminant}
\externaldocument[derham-]{derham}
\externaldocument[local-cohomology-]{local-cohomology}
\externaldocument[algebraization-]{algebraization}
\externaldocument[curves-]{curves}
\externaldocument[resolve-]{resolve}
\externaldocument[models-]{models}
\externaldocument[functors-]{functors}
\externaldocument[equiv-]{equiv}
\externaldocument[pione-]{pione}
\externaldocument[etale-cohomology-]{etale-cohomology}
\externaldocument[proetale-]{proetale}
\externaldocument[relative-cycles-]{relative-cycles}
\externaldocument[more-etale-]{more-etale}
\externaldocument[trace-]{trace}
\externaldocument[crystalline-]{crystalline}
\externaldocument[spaces-]{spaces}
\externaldocument[spaces-properties-]{spaces-properties}
\externaldocument[spaces-morphisms-]{spaces-morphisms}
\externaldocument[decent-spaces-]{decent-spaces}
\externaldocument[spaces-cohomology-]{spaces-cohomology}
\externaldocument[spaces-limits-]{spaces-limits}
\externaldocument[spaces-divisors-]{spaces-divisors}
\externaldocument[spaces-over-fields-]{spaces-over-fields}
\externaldocument[spaces-topologies-]{spaces-topologies}
\externaldocument[spaces-descent-]{spaces-descent}
\externaldocument[spaces-perfect-]{spaces-perfect}
\externaldocument[spaces-more-morphisms-]{spaces-more-morphisms}
\externaldocument[spaces-flat-]{spaces-flat}
\externaldocument[spaces-groupoids-]{spaces-groupoids}
\externaldocument[spaces-more-groupoids-]{spaces-more-groupoids}
\externaldocument[bootstrap-]{bootstrap}
\externaldocument[spaces-pushouts-]{spaces-pushouts}
\externaldocument[spaces-chow-]{spaces-chow}
\externaldocument[groupoids-quotients-]{groupoids-quotients}
\externaldocument[spaces-more-cohomology-]{spaces-more-cohomology}
\externaldocument[spaces-simplicial-]{spaces-simplicial}
\externaldocument[spaces-duality-]{spaces-duality}
\externaldocument[formal-spaces-]{formal-spaces}
\externaldocument[restricted-]{restricted}
\externaldocument[spaces-resolve-]{spaces-resolve}
\externaldocument[formal-defos-]{formal-defos}
\externaldocument[defos-]{defos}
\externaldocument[cotangent-]{cotangent}
\externaldocument[examples-defos-]{examples-defos}
\externaldocument[algebraic-]{algebraic}
\externaldocument[examples-stacks-]{examples-stacks}
\externaldocument[stacks-sheaves-]{stacks-sheaves}
\externaldocument[criteria-]{criteria}
\externaldocument[artin-]{artin}
\externaldocument[quot-]{quot}
\externaldocument[stacks-properties-]{stacks-properties}
\externaldocument[stacks-morphisms-]{stacks-morphisms}
\externaldocument[stacks-limits-]{stacks-limits}
\externaldocument[stacks-cohomology-]{stacks-cohomology}
\externaldocument[stacks-perfect-]{stacks-perfect}
\externaldocument[stacks-introduction-]{stacks-introduction}
\externaldocument[stacks-more-morphisms-]{stacks-more-morphisms}
\externaldocument[stacks-geometry-]{stacks-geometry}
\externaldocument[moduli-]{moduli}
\externaldocument[moduli-curves-]{moduli-curves}
\externaldocument[examples-]{examples}
\externaldocument[exercises-]{exercises}
\externaldocument[guide-]{guide}
\externaldocument[desirables-]{desirables}
\externaldocument[coding-]{coding}
\externaldocument[obsolete-]{obsolete}
\externaldocument[fdl-]{fdl}
\externaldocument[index-]{index}

% Theorem environments.
%
\theoremstyle{plain}
\newtheorem{theorem}[subsection]{Theorem}
\newtheorem{proposition}[subsection]{Proposition}
\newtheorem{lemma}[subsection]{Lemma}

\theoremstyle{definition}
\newtheorem{definition}[subsection]{Definition}
\newtheorem{example}[subsection]{Example}
\newtheorem{exercise}[subsection]{Exercise}
\newtheorem{situation}[subsection]{Situation}

\theoremstyle{remark}
\newtheorem{remark}[subsection]{Remark}
\newtheorem{remarks}[subsection]{Remarks}

\numberwithin{equation}{subsection}

% Macros
%
\def\lim{\mathop{\mathrm{lim}}\nolimits}
\def\colim{\mathop{\mathrm{colim}}\nolimits}
\def\Spec{\mathop{\mathrm{Spec}}}
\def\Hom{\mathop{\mathrm{Hom}}\nolimits}
\def\Ext{\mathop{\mathrm{Ext}}\nolimits}
\def\SheafHom{\mathop{\mathcal{H}\!\mathit{om}}\nolimits}
\def\SheafExt{\mathop{\mathcal{E}\!\mathit{xt}}\nolimits}
\def\Sch{\mathit{Sch}}
\def\Mor{\mathop{\mathrm{Mor}}\nolimits}
\def\Ob{\mathop{\mathrm{Ob}}\nolimits}
\def\Sh{\mathop{\mathit{Sh}}\nolimits}
\def\NL{\mathop{N\!L}\nolimits}
\def\CH{\mathop{\mathrm{CH}}\nolimits}
\def\proetale{{pro\text{-}\acute{e}tale}}
\def\etale{{\acute{e}tale}}
\def\QCoh{\mathit{QCoh}}
\def\Ker{\mathop{\mathrm{Ker}}}
\def\Im{\mathop{\mathrm{Im}}}
\def\Coker{\mathop{\mathrm{Coker}}}
\def\Coim{\mathop{\mathrm{Coim}}}

% Boxtimes
%
\DeclareMathSymbol{\boxtimes}{\mathbin}{AMSa}{"02}

%
% Macros for moduli stacks/spaces
%
\def\QCohstack{\mathcal{QC}\!\mathit{oh}}
\def\Cohstack{\mathcal{C}\!\mathit{oh}}
\def\Spacesstack{\mathcal{S}\!\mathit{paces}}
\def\Quotfunctor{\mathrm{Quot}}
\def\Hilbfunctor{\mathrm{Hilb}}
\def\Curvesstack{\mathcal{C}\!\mathit{urves}}
\def\Polarizedstack{\mathcal{P}\!\mathit{olarized}}
\def\Complexesstack{\mathcal{C}\!\mathit{omplexes}}
% \Pic is the operator that assigns to X its picard group, usage \Pic(X)
% \Picardstack_{X/B} denotes the Picard stack of X over B
% \Picardfunctor_{X/B} denotes the Picard functor of X over B
\def\Pic{\mathop{\mathrm{Pic}}\nolimits}
\def\Picardstack{\mathcal{P}\!\mathit{ic}}
\def\Picardfunctor{\mathrm{Pic}}
\def\Deformationcategory{\mathcal{D}\!\mathit{ef}}
\usepackage{cleveref}
\crefname{originallemma}{original lemma}{original lemmas}
\Crefname{originallemma}{Original lemma}{Original lemmas}
\makeatletter
\newlabel{lemma-compare-constructions}{{01I7}{source}{}{sourcepointer-1}{}}
\newlabel{lemma-compare-constructions@cref}{{[originallemma][0][]01I7}{[1][0][]source}{}{}{}}
\makeatother
\crefname{originallemma}{original lemma}{original lemmas}
\Crefname{originallemma}{Original lemma}{Original lemmas}
\makeatletter
\newlabel{lemma-spec-sheaves}{{01HV}{source}{}{sourcepointer-2}{}}
\newlabel{lemma-spec-sheaves@cref}{{[originallemma][0][]01HV}{[1][0][]source}{}{}{}}
\makeatother
\title{Two out of three for quasi-coherence in a short exact sequence}
\author{The Stacks Project Authors}
\date{stacks-project:schemes.tex:local-snapshot-20261012; GNU FDL 1.2 or later}
\begin{document}
\maketitle
\noindent\textit{Editorial note.} Counted span is the author's own proof environment, source lines 1689-1782, following the statement at 1674-1687.\par\medskip
\noindent\textit{Editorial note.} The author flags the alternative H\textasciicircum{}1 = 0 route as postponed and uses the ideal trick instead; the trick is carried out fully explicitly.\par\medskip
\noindent\textit{Allowed standard prerequisites (editorial).}\begin{itemize}
\item The tilde functor is exact. lemma-spec-sheaves (same chapter, cited at \textasciitilde{}1690-1692). Engine for the two easy cases of the trichotomy.
\item Mapping property and sections of tilde sheaves. lemma-compare-constructions (same chapter, cited at \textasciitilde{}1701-1702, \textasciitilde{}1750, \textasciitilde{}1775). Builds the comparison diagram and identifies F\_1 sections with (M\_1)-localizations.
\item Snake lemma. Standard homological algebra (used at \textasciitilde{}1718). Concludes quasi-coherence of F\_2 once 0 -> M\_1 -> M\_2 -> M\_3 -> 0 is exact.
\end{itemize}\medskip
\ifdefined\hypertarget\hypertarget{sourcepointer-1}{}\fi
\noindent\textit{Source reference.} Original lemma 01I7: sources/stacks-schemes/pinned-original-source.tex, source lines 1295--1295. Navigation pointer only: the counted proof cites this statement by \textbackslash{}ref as a standing input; the deduction the counted proof has to make is carried out inside the counted span, and no part of that deduction is delegated to this pointer.\par\medskip
\ifdefined\hypertarget\hypertarget{sourcepointer-2}{}\fi
\noindent\textit{Source reference.} Original lemma 01HV: sources/stacks-schemes/pinned-original-source.tex, source lines 909--909. Navigation pointer only: the counted proof cites this statement by \textbackslash{}ref as a standing input; the deduction the counted proof has to make is carried out inside the counted span, and no part of that deduction is delegated to this pointer.\par\medskip
\begin{lemma}
\label{lemma-extension-quasi-coherent}
Let $(X, \mathcal{O}_X) = (\Spec(R), \mathcal{O}_{\Spec(R)})$
be an affine scheme. Suppose that
$$
0 \to
\mathcal{F}_1 \to
\mathcal{F}_2 \to
\mathcal{F}_3 \to
0
$$
is a short exact sequence of sheaves of $\mathcal{O}_X$-modules.
If two out of three are quasi-coherent then so is the third.
\end{lemma}

\begin{proof}
This is clear in case both $\mathcal{F}_1$ and $\mathcal{F}_2$ are
quasi-coherent because the functor $M \mapsto \widetilde M$
is exact, see Lemma \ref{lemma-spec-sheaves}.
Similarly in case both $\mathcal{F}_2$ and $\mathcal{F}_3$ are
quasi-coherent. Now, suppose that $\mathcal{F}_1 = \widetilde M_1$ and
$\mathcal{F}_3 = \widetilde M_3$ are quasi-coherent.
Set $M_2 = \Gamma(X, \mathcal{F}_2)$. We claim it suffices to show that
the sequence
$$
0 \to M_1 \to M_2 \to M_3 \to 0
$$
is exact. Namely, if this is the case, then (by using the mapping
property of Lemma \ref{lemma-compare-constructions}) we get a commutative
diagram
$$
\xymatrix{
0 \ar[r] &
\widetilde M_1 \ar[r] \ar[d] &
\widetilde M_2 \ar[r] \ar[d] &
\widetilde M_3 \ar[r] \ar[d] &
0 \\
0 \ar[r] &
\mathcal{F}_1 \ar[r] &
\mathcal{F}_2 \ar[r] &
\mathcal{F}_3 \ar[r] &
0
}
$$
and we win by the snake lemma.

\medskip\noindent
The ``correct'' argument here would be to show first
that $H^1(X, \mathcal{F}) = 0$ for any quasi-coherent sheaf $\mathcal{F}$.
This is actually not all that hard, but it is perhaps better to postpone
this till later. Instead we use a small trick.

\medskip\noindent
Pick $m \in M_3 = \Gamma(X, \mathcal{F}_3)$.
Consider the following set
$$
I = \{ f \in R \mid \text{the element }fm\text{ comes from }M_2\}.
$$
Clearly this is an ideal. It suffices to show $1 \in I$.
Hence it suffices to show that for any prime $\mathfrak p$
there exists an $f \in I$, $f \not\in \mathfrak p$.
Let $x \in X$ be the point corresponding to $\mathfrak p$.
Because surjectivity can be checked on stalks
there exists an open neighbourhood $U$ of $x$ such that
$m|_U$ comes from a local section $s \in \mathcal{F}_2(U)$.
In fact we may assume that $U = D(f)$ is a standard open,
i.e., $f \in R$, $f \not \in \mathfrak p$. We will show
that for some $N \gg 0$ we have $f^N \in I$, which
will finish the proof.

\medskip\noindent
Take any point $z \in V(f)$, say corresponding to the
prime $\mathfrak q \subset R$. We can also find a $g \in R$,
$g \not \in \mathfrak q$ such that $m|_{D(g)}$ lifts
to some $s' \in \mathcal{F}_2(D(g))$.
Consider the difference $s|_{D(fg)} - s'|_{D(fg)}$.
This is an element $m'$ of $\mathcal{F}_1(D(fg)) = (M_1)_{fg}$.
For some integer $n = n(z)$ the element $f^n m'$ comes
from some $m'_1 \in (M_1)_g$. We see that
$f^n s$ extends to a section $\sigma$ of $\mathcal{F}_2$ on $D(f) \cup D(g)$
because it agrees with the restriction of
$f^n s' + m'_1$ on $D(f) \cap D(g) = D(fg)$.
Moreover, $\sigma$ maps to the restriction of $f^n m$
to $D(f) \cup D(g)$.

\medskip\noindent
Since $V(f)$ is quasi-compact, there exists a finite list
of elements $g_1, \ldots, g_m \in R$ such that
$V(f) \subset \bigcup D(g_j)$, an integer $n > 0$ and sections
$\sigma_j \in \mathcal{F}_2(D(f) \cup D(g_j))$ such that
$\sigma_j|_{D(f)} = f^n s$ and $\sigma_j$ maps to the section
$f^nm|_{D(f) \cup D(g_j)}$ of $\mathcal{F}_3$.
Consider the differences
$$
\sigma_j|_{D(f) \cup D(g_jg_k)}
-
\sigma_k|_{D(f) \cup D(g_jg_k)}.
$$
These correspond to sections of $\mathcal{F}_1$
over $D(f) \cup D(g_jg_k)$ which are zero
on $D(f)$. In particular their images in
$\mathcal{F}_1(D(g_jg_k)) = (M_1)_{g_jg_k}$
are zero in $(M_1)_{g_jg_kf}$.
Thus some high power of $f$ kills each and every one of these.
In other words, the elements $f^N \sigma_j$, for some $N \gg 0$
satisfy the glueing condition of the sheaf property and
give rise to a section $\sigma $ of $\mathcal{F}_2$
over $\bigcup (D(f) \cup D(g_j)) = X$ as desired.
\end{proof}

\end{document}
